

\textit{Bearbeitet von Mohammad Badr Eddin}\\

Für die Positionsregelung des Greifers ist eine zuverlässige Rückführung
der Position der Motorwelle erforderlich. Hierfür wird ein magnetischer
Positionssensor eingesetzt, der die Drehstellung der Welle berührungslos
erfasst. Das erfasste Feedback wird anschließend an den Mikrocontroller
übertragen und für die Positionsregelung verwendet.

\section{Positionsrückführung mit dem AS5600}
\label{sec:encoder}

Für die Erfassung der Motorposition wird ein magnetischer
Drehpositionssensor des Typs AS5600 eingesetzt. Der Sensor arbeitet
berührungslos und bestimmt die Winkelstellung eines Permanentmagneten,
der auf der Motorwelle angebracht ist. Dadurch kann die Position der
Motorwelle erfasst werden, ohne dass eine mechanische Verbindung
zwischen Sensor und Magnet erforderlich ist.

Der AS5600 besitzt eine interne 12-Bit-Auswertung und kann die
Winkelstellung über einen Bereich von $0^\circ$ bis $360^\circ$
erfassen. Damit eignet sich der Sensor für die Rückführung der
Motorposition innerhalb des Greiferantriebs. Der Sensor wird
nicht direkt zur Messung der linearen Fingerposition verwendet.
Stattdessen wird zunächst die Drehbewegung der Motorwelle erfasst und
die daraus resultierende lineare Bewegung des Greifers später über die
Kinematik des Zahnstangen-Ritzel-Systems bestimmt.


\subsection{Funktionsprinzip des magnetischen Sensors}

Die Winkelmessung des AS5600 basiert auf dem Hall-Effekt. Hierzu wird
ein Permanentmagnet so angeordnet, dass sich das Magnetfeld im Bereich
des Sensors mit der Drehung der Motorwelle verändert. Der AS5600
erfasst dieses Magnetfeld berührungslos und bestimmt daraus die
Winkelstellung des Magneten.

Für die vorgesehene Anwendung wird ein diametral magnetisierter
Permanentmagnet verwendet. Bei einer Drehung der Motorwelle dreht sich
der Magnet entsprechend mit. Die dadurch entstehende Änderung der
Magnetfeldrichtung wird vom Sensor erfasst und intern in einen
Winkelwert umgerechnet.

Die interne Auflösung des Sensors beträgt 12 Bit. Der vollständige
Messbereich wird dadurch in $4096$ diskrete Positionen unterteilt. Die
Winkelauflösung pro Messschritt beträgt somit

\begin{equation}
\Delta\varphi
=
\frac{360^\circ}{4096}
\approx 0{,}0879^\circ.
\end{equation}

Der ausgegebene Winkelwert stellt somit eine absolute Position
innerhalb einer Umdrehung dar. Eine Erfassung mehrerer Umdrehungen
erfolgt nicht direkt durch den Sensor, sondern kann bei Bedarf durch
eine entsprechende Verarbeitung der aufeinanderfolgenden Messwerte
in der Firmware realisiert werden.

\subsection{Mechanische Integration in den Greifer}
\label{sec:encoder_montage}

Der AS5600 ist an der Rückseite des Antriebs angeordnet. Der
Permanentmagnet wurde auf der hinteren Motorwelle befestigt und dreht
sich somit gemeinsam mit der Welle. Zur Befestigung wurde der Magnet
mit Schnellkleber auf der Rückseite der Welle fixiert. Der AS5600
selbst wurde über eine separate Halterung gegenüber dem Magneten
befestigt.

Für eine zuverlässige Winkelmessung ist die Positionierung des
Magneten gegenüber dem Sensor entscheidend. Der Magnet muss
korrekt zum Sensor ausgerichtet sein und sich im vorgesehenen
Messbereich des AS5600 befinden. Die Halterung sorgt dafür, dass der Sensor gegenüber dem Magneten in der vorgesehenen Position bleibt.


% \textit{Bearbeitet von Mohammad Badr Eddin}\\

% Für die Positionsregelung des Greifers ist eine zuverlässige Rückführung der Motorwelleposition erforderlich. Dazu wird die aktuelle Stellung des Antriebs über einen magnetischen Positionssensor erfasst und dem Mikrocontroller zur Verarbeitung bereitgestellt.
% Die erfasste Feedback ist in dem Fall eine Eingangsgröße für die Positionsregelung. 

% \section{Positionsrückführung mit dem AS5600}
% \label{sec:encoder}

% Für die Erfassung der Motorposition wird ein magnetischer
% Drehpositionssensor des Typs AS5600 eingesetzt. Der Sensor arbeitet
% kontaktlos und bestimmt die Winkelposition eines diametral magnetisierten
% Permanentmagneten. Die interne Auswertung gibt eine absolute
% Winkelposition mit einer Auflösung von 12 Bit über einen Bereich von
% $0^\circ$ bis $360^\circ$ aus.

% Der Sensor ist über die I\textsuperscript{2}C-Schnittstelle mit dem
% STM32F767 verbunden. Der AS5600 besitzt die 7-Bit-I\textsuperscript{2}C-
% Adresse $0x36$. Für die Positionsrückführung werden die
% Register \texttt{RAW\_ANGLE} bei den Adressen $0x0C$ und $0x0D$ verwendet.
% Diese liefern den 12-Bit-Rohwert der gemessenen Winkelposition. Zusätzlich
% wird das Statusregister bei Adresse $0x0B$ zur Erkennung des Magneten
% ausgewertet \cite{AMS}.

% \subsection{Auslesen des Encoderwertes}

% Die Firmware initialisiert den Encoder durch eine I\textsuperscript{2}C-
% Geräteerkennung. Anschließend wird der Rohwert zyklisch über einen
% Speicher-Lesezugriff aus den beiden \texttt{RAW\_ANGLE}-Registern
% ausgelesen. Die beiden Bytes werden zu einem 12-Bit-Wert zusammengesetzt:

% \begin{equation}
% N_{\mathrm{enc}}
% =
% \left(N_{\mathrm{high}} \ll 8\right)
% \mathbin{|}
% N_{\mathrm{low}}
% \end{equation}

% Anschließend werden die nicht verwendeten oberen Bits maskiert. Der
% resultierende Wertebereich beträgt somit

% \begin{equation}
% N_{\mathrm{enc}} \in [0,4095].
% \end{equation}

% Die Winkelauflösung des Sensors beträgt damit

% \begin{equation}
% \Delta\varphi
% =
% \frac{360^\circ}{4096}
% \approx 0{,}0879^\circ.
% \end{equation}


% \subsection{Umrechnung in die lineare Greiferposition}

% Für die Regelung wird nicht der Winkel des Encoders, sondern die daraus
% resultierende lineare Position des Greifers verwendet. Das Ritzel besitzt
% einen Teilkreisdurchmesser von $d=17\,\si{mm}$. Eine vollständige
% Umdrehung entspricht daher einer Umfangsstrecke von

% \begin{equation}
% s_{\mathrm{rev}}
% =
% \pi d
% =
% \pi \cdot 17\,\si{mm}
% \approx 53{,}41\,\si{mm}.
% \end{equation}

% Unter der Annahme einer direkten Kopplung zwischen Encoderwelle und Ritzel
% ergibt sich daraus die lineare Auflösung

% \begin{equation}
% \Delta x_{\mathrm{enc}}
% =
% \frac{s_{\mathrm{rev}}}{4096}
% =
% \frac{53{,}41\,\si{mm}}{4096}
% \approx 0{,}0130\,\si{mm}.
% \end{equation}

% Ein Encoder-Count entspricht somit etwa $\SI{13}{\micro\meter}$ linearer
% Bewegung am Ritzel.
% % +++++++++++++++++++++++++++++++++++++++++++++++++++++++++++++++++
% % Sensorik
% % +++++++++++++++++++++++++++++++++++++++++++++++++++++++++++++++++
% \section{Sensorik zur Strommessung}
% \textit{[Begründung Auswahl und Auswertung geeigneter Sensorik zur Strommessung]}


% % +++++++++++++++++++++++++++++++++++++++++++++++++++++++++++++++++
% % Signalaufbereitung
% % +++++++++++++++++++++++++++++++++++++++++++++++++++++++++++++++++
% \section{Signalaufbereitung}
% % Filterung, Skalierung der Rohdaten für die Regelung.
% \textit{[Filterung und Skalierung der Rohdaten für die Regelung]}